%% fuer Drucker "bilby":
%% dvips -O 5mm,1mm

%% fuer Drucker "goofy":
%% dvips600 -O 0mm,10mm Magdeburg

\documentclass[portrait]{seminar}
\usepackage{german,colordvi,boldmath,semcolor,semlayer,rotate,epsf,fancybox,mathsym}
%\input{farbunilogo}
%\input{unilogo}
\usepackage{amssymb}
\UseLegacyTextSymbols

\def\dummy#1{}
\overlaystrue
\def\overlay#1{{}}

\slidewidth9.5in
\slideheight6.7in
%\def\slidetopmargin{\iflandscape.2in\else0.8in\fi}
\def\slidetopmargin{\iflandscape.3in\else0.8in\fi}
\def\slidebottommargin{\iflandscape1.0in\else0.4in\fi}
\def\slideleftmargin{\iflandscape.2in\else.5in\fi}
\def\sliderightmargin{\iflandscape1.0in\else.8in\fi}

\newcommand{\cH}{{\cal H}}

\slidesmag{4}
\centerslidesfalse
\slideframewidth0.5pt
%\slideframe{none}

%\portraitonly
%\landscapeonly
%\def\leftsliderotation#1{\rotate[l]{#1}}
%\sliderotation{left}
\rotateheaderstrue

%%%%%%%%%%%%%%%%%%%%%%%%%%%%%%%%%%%%%%%%%%%%%%
%% Makros
%%%%%%%%%%%%%%%%%%%%%%%%%%%%%%%%%%%%%%%%%%%%%%
\def\bra#1{\left<#1\right|}
\def\ket#1{\left|#1\right>}

\def\trace{\mathop{\rm tr}\nolimits}

\newcommand{\heading}[1]{%
  \begin{center}
    \Large
    \bfseries
    \shadowbox{#1}%
  \end{center}
  \vspace{1ex minus 1ex}%
}
\newbox{\logobox}
%\setbox\logobox=\hbox{\raise 0.17\baselineskip\hbox{\epsfbox{logo.eps}}}

\def\logo{\usebox{\logobox}}
\newpagestyle{MH}{\hss\hbox to \hsize{PhysComp96\hss\thepage}\hss}%
{\hss\hbox to \hsize{${\cal IAKS}$\ \logo\ UniKA\hfil}\hss}

\newpagestyle{nohead}{}{\hss\hbox to \hsize{\tiny IAKS
\ \logo\ Universit"at Karlsruhe\hfil Prof.~Dr.~Th.~Beth}\hss}

\newpagestyle{nohead_beth}{}{\hss\hbox to \hsize{\tiny${\cal IAKS}$\ \logo\
UniKA\hfil}\hss}

\newcommand{\tl}{\tilde}
\newcommand{\onemat}{{\mathbf 1}}

\begin{document}

\begin{slide}
{$ $}
\vfill
\begin{center}
{\Large \blue Computational Power of Hamiltonians\\ in Quantum Computing} 

\vfill Pawel Wocjan

{\small Institut f"ur Algorithmen und Kognitive Systeme}

{\small Universit"at Karlsruhe}
\end{center}
\vfill
\end{slide}

\begin{slide}
{\large {\bf Quanteninformationstheorie}}
\vfill
\begin{itemize}
\item Quantenkomlexit\"atstheorie
\item Quantenkryptographie
\item Quantenkommunikation
\end{itemize}
\vfill
\end{slide}

\begin{slide}
{\large {\bf Fast control limit}}
\vfill
strength of the control Hamiltonians arbitrarily large $\Longrightarrow$

{\blue control unitaries} arbitrarily fast
\vfill
{\red Complexity of a unitary $U$}: minimal time
\[
U=\exp(-i H t_N) V_N \cdots \exp(-i H t_2) V_2 \exp(-i H t_1) V_1
\]
\vfill
express it as the evolution according to the piecewise constant Hamiltonian
\[
U = U_N \exp(-i U_N^\dagger H U_N t_N) \cdots \exp(-i U_2^\dagger H U_2 t_2)
    \exp(-i U_1^\dagger H U_1 t_1)
\]
\vfill
\end{slide}


\begin{slide}
{\large {\bf Average Hamiltonian Theory}}
\vfill
{\blue Dyson expansion}
\[
U(t)={\cal T} \exp(\int_{\tau=0}^t d\tau H(\tau))=\exp(-i\bar{H}t)
\]
\vfill
{\blue Magnus expansion} $\bar{H}=\bar{H}^{(0)}+\bar{H}^{(1)}+\cdots$
\begin{eqnarray*}
{\red\bar{H}^{(0)}} & {\red =} & {\red \frac{1}{t}\int_0^t d\tau
H(\tau)\,,\quad\quad\quad\quad\ \quad\quad\quad\quad
\|\bar{H}^{(0)}\|\le 1} \\
\bar{H}^{(1)} & = & \frac{-i}{2t}\int_0^t
d\tau' \int_0^t\tau'' [H(\tau'),H(\tau'')]\,,\quad
\|\bar{H}^{(1)}\|\le t/2
\end{eqnarray*}
\vfill
{\red $\Longrightarrow$ $U(t)$ essentially determined by $\bar{H}^{(0)}$ for
small $t$}
\vfill
\end{slide}


\begin{slide}
{\large {\bf Simulating Hamiltonians/First order simulation}}
\vfill
{\blue $H$ can simulate $\tilde{H}$ with overhead $1$}, 
written $\tilde{H}\prec H$, iff 
\[
\tilde{H}=\sum_{j=1}^N p_j U_j^\dagger H U_j
\]
with $\sum_j p_j = 1$ and $U_j$ control operations, i.e.

{\red $\tilde{H}$ is a convex sum of unitary conjugates of $H$}
\vfill
$H$ can simulate $\tilde{H}$ with overhead {\red $\tau$} iff 
$\tilde{H} \prec {\red \tau} H$ 
\vfill 
{\red time-optimal simulation} problem $\longrightarrow$
{\red convex optimization} problem 
\vfill
\end{slide}

\begin{slide}
{\large {\bf Quantum Networks}}
\vfill
system consisting of {\blue $n$ nodes}  
$\quad {\cal H}=\C^d\otimes\C^d\otimes\cdots\otimes\C^d$ 
\vfill
{\blue pair-interaction Hamiltonian}
\[
H = \sum_{k<l} \sum_{\alpha\beta} J_{kl;\alpha\beta} 
\sigma_\alpha^k\sigma_\beta^l\,,\quad \sigma_\alpha \mbox{ ONB of } su(d)
\]
{\blue coupling matrix}
\[
J=\left(
\begin{array}{c|c|c|c|c}
0      & J_{12} & J_{13} & \cdots & J_{1n} \\ \hline
J_{21} & 0      & J_{23} & \cdots & J_{2n} \\ \hline
\vdots & \vdots &        &        &        \\ \hline
J_{n1} & J_{n2} & J_{n3} &        & 0 
\end{array}
\right)\in\R^{(d^2-1)\,n\times (d^2-1)\,n}
\]
{\blue control operations} are all {\red local} operations in
\[
SU(d)\otimes SU(d)\otimes\cdots\otimes SU(d)
\]
\end{slide}

\begin{slide}
{\bf Decoupling/Switching off the time evolution} 
\vfill 
Conjugating by unitaries of an {\blue error basis} ${\cal E}=\{U_1={\mathbf
1},U_2,\ldots,U_{d^2}\}$
\[
\sum_j U_j^\dagger\, a\, U_j = {\mathbf 0}\quad\quad
\mbox{for all } a\in su(d)
\]
cancels all matrices in $su(d)$
\vfill
Combinatorial problem: which sequences of conjugations allow to switch off
all interactions?
\vfill
Efficient solution: decoupling in networks based on
{\blue orthogonal arrays} 

number of conjugations grows linearly with $n$
\vfill
\end{slide}

\begin{slide}
{\large {\bf Action of control operations on the coupling matrix $J$}}
\vfill
\begin{center}
{\blue Conjugation of the pair-interaction Hamiltonian $H$ by control
operations}

corresponds to

{\blue conjugation of $J$ by a direct sum of orthogonal matrices}
\end{center}
\[
\begin{array}{ccc}
H & \mapsto & (u_1\otimes u_2\otimes\cdots\otimes u_n)^\dagger\, H\, 
(u_1\otimes u_2\otimes\cdots\otimes u_n) \\
J & \mapsto & (U_1\oplus U_2\oplus\cdots\oplus U_n)\, J\,
(U_1\oplus U_2\oplus\cdots\oplus U_n)^T
\end{array}
\]
\vfill
\end{slide}

\begin{slide}
{\large {\bf Lower bounds on time overhead}}
\vfill
If $H$ can simulate $\tilde{H}$ with overhead $\tau$ then
\[
\tilde{J}/\tau = \sum_{j=1}^N p_j 
(U_{j1}\oplus U_{j2}\oplus\cdots\oplus U_{jn})
\,J\, (U_{j1}\oplus U_{j2}\oplus\cdots\oplus U_{jn})^T
\]
convex sum of conjugates of $J$
\vfill
$\Rightarrow$ necessary condition given by {\blue majorization criterion}:\\
\begin{center}
the spectrum Spec($\tilde{J}$) majorizes the spectrum Spec($\tau J$)
\end{center}
\vfill
\end{slide}


\begin{slide}
{\large {\bf Application I: Time reversal}}
\vfill
Time reversal: $H \mapsto -H$ (refocusing, spin echoes)
\vfill
dipole-dipole coupling between all nodes (easy to invert)
\[
H=\sum_{k<l} w_{kl} (3\,\sigma_z\otimes\sigma_z - \sum_{\alpha=x,y,z} 
\sigma_\alpha\otimes\sigma_\alpha)\,,\quad w_{kl}\neq 0
\]
\vfill
lower bound on time overhead is $2$
\vfill
this bound is tight (well known-scheme in NMR)
\vfill
\end{slide}


\begin{slide}
{\large {\bf Application I: Time reversal}}
\vfill
strong scalar coupling between all nodes (difficult to invert)
\[
H = \sum_{k<l} w_{kl}
(\sigma_x^k\sigma_x^l + \sigma_y^k\sigma_y^l + \sigma_z^k\sigma_z^l)\,,\quad
w_{kl}\neq 0
\]
\vfill
lower bound on the overhead $n-1$

the inverted time-evolution is necessarily slower by the factor $(n-1)$
\vfill
upper bound $cn-1$ (decoupling scheme based on orthogonal arrays)
\vfill
\end{slide}

\begin{slide}
{\large {\bf Application II: Hierarchy of Hamiltonians}}
\vfill
\[
H_1 = \sum_{kl} w_{kl} 
(\sigma^k_x\otimes\sigma^l_x\, {\red +}\, \sigma^k_y\otimes\sigma^l_y)
\]
\[
H_2 = \sum_{kl} w_{kl} 
(\sigma^k_x\otimes\sigma^l_x\, {\blue -}\, \sigma^k_y\otimes\sigma^l_y)
\]
\vfill
$H_1$ cannot simulate $H_2$ with overhead less than $(n-1)$

$H_2$ can simulate $H_1$ with overhead $2$
\vfill
\end{slide}

\begin{slide}
{\large {\bf Application III: Switching off {\red some} interactions}}
\vfill
Combinatorial problem: cancel some interactions
\vfill
Are the other interactions necessarily disturbed?
\begin{figure}
\centerline{
\epsfbox[0 0 285 98]{HonnefGraphen.eps}}
\end{figure}
\vfill 
Yes, for some graphs $\rightarrow$ lower bound on time overhead
given by eigenvalues of the adjacency matrix \vfill

Upper bound given by chromatic index of interaction graph
\end{slide}

\begin{slide}
{\bf Optimal simulation problem equivalent to separability problem}
\vfill
complete $zz$-Hamiltonian $H$: interaction between all nodes are of the
form $\sigma_z\otimes \sigma_z$
\vfill
A Hamiltonian with coupling matrix $J$ can be simulated with overhead 
$1$ iff there is a separable $n$-qubit state $\rho$ such that $J$ is the
$3n\times 3n$-matrix
\[
\left( tr(\rho\sigma_\alpha^k\sigma_l^l)\right)_{kl\alpha\beta}
\]
\vfill
\begin{center}
{\blue maximal strength of pair-interactions implementable with overhead $1$}

corresponds to

{\blue maximal strength of pair-correlations in a separable state}
\end{center}
\vfill
\end{slide}

\begin{slide}
{\large {\bf Universal simulation using finite control groups}}
\vfill
{\blue Universal simulation} does not require all operations in $SU(d)$

Certain subgroups are sufficient; called {\blue transformer groups}
\vfill 
{\blue Finite transformer groups are characterized in terms of group 
characters} 
\vfill
\end{slide}

\begin{slide}
{\bf {\large Conclusions}}
\begin{itemize}
\item Upper and lower bounds on simulation complexity are given by
graph-theoretical notions like chromatic index, chromatic number and
graph spectra
\item Efficient decoupling schemes can be constructed by using 
orthogonal arrays and error basis
\item Universal simulation does not require the ability to perform all
local unitary operations; we have found finite control groups that are
sufficient
\end{itemize}
{\tiny
{\it Simulating Arbitrary Pair-interactions by a Given Hamiltonian: 
Graph-theoretical Bounds on the Time Complexity},
quant-ph/0106077, to appear in QIC

{\it Simulating Hamiltonians in Quantum Networks: Efficient schemes and
Complexity Bounds}, quant-ph/0109088, to appear in Phys. Rev. A

{\it Universal Simulation of Hamiltonians Using a Finite Set of Control 
Operations}, quant-ph/0109063, to appear in QIC

{\it Complexity of inverting n-spin interactions: Arrow of time in quantum 
control}, quant-ph/0106085

{\it Lower Bound on the Chromatic Number by Spectra of Weighted Adjacency 
Matrices}, cs.DM/0112023
}
\end{slide}
\end{document}









